\chapter{Lecture 22: Tranformation by the cdf, and expected values}

\section{Transformation by the cdf} \label{U}
In lecture 21 we learned how to transform a random variable, $X$, and we learned some tools for finding the distribution of $Y$ when $Y = g(X)$. We looked at cases when the transformation, $g$, was invertible, i.e., strictly increasing or decreasing (where method of distributions applies), and we looked at cases where $g$ was non-invertible (where we needed to use Leibniz's rule to evaluate a probability integral). \\

\noindent Now we will learn a particularly interesting and useful application of transformations. In particular, we will look at the transformation $\boldsymbol{U = F_X(X)}$. Here, we have transformed a random variable, $X$, by using its own cdf, $F_X$, as the transformation. So, normally for a transformation we would write $Y = g(X)$, but in this case we have replaced $g$ with $F_X$, and we have used the symbol $U$ instead of $Y$. \\

\noindent Why are we writing $U$ instead of $Y$? As we will soon see, it is because \textbf{every time we transform $\boldsymbol{X}$ using its own cdf, $\boldsymbol{F_X,}$ the resulting variable will have the U($\boldsymbol{0,1}$) distribution}; i.e., it will be uniformly distributed on $[0,1];$ we use $U$ to represent a continuous random variable whose distribution is U($0,1$). \\

\noindent So, this equation, $U = F_X(X)$, tells us which transformation to use (i.e., $F_X$) to transform any random variable, $X$, into $U$.  \\

\noindent This is an interesting and useful fact that relates the uniform distribution to other distributions. \\

\noindent Suppose, instead, that we want to start with a uniformly distributed random variable, $U$ (something which is comparatively easy to generate on a computer), and we wish to find the right function to transform $U$ into some continuous random variable, $X$, where $X$ is distributed however we want (i.e., we want to choose the cdf, $F_X$, ourselves). \\

\noindent Solving for $X$ in the equation $U = F_X(X)$, gives $\boldsymbol{X = F^{-1}_X(U)}$. This equation tells us that we can choose any cdf, $F_X$, find $F^{-1}_X$, then apply this transformation to $U$, to obtain a random variable, $X$, whose cdf is $F_X$. This is a handy result, since many computer packages can easily be used to generate uniformly distributed `random' numbers. By applying $F^{-1}_X$ to them, we can generate new random numbers which follow the distribution whose cdf is $F_X$. This is a useful way to generate `random' numbers following whatever distribution that we want. \\

\noindent Note that for any continuous cdf, $F_X$, there always exist some interval, $[a,b]$, on which $F_X$ is strictly increasing (see section \ref{cdf(continuous)}). Also, $F_X$ is flat everwhere else. Hence we need to restrict the domain of $F_X$ to $[a,b]$, when finding $F^{-1}(x)$. We spoke about this concept in more detail in section \ref{cdf strictly increasing 1} when we discussed inverse functions and percentiles.
%\begin{center}
%\begin{tikzpicture}
%\axesmacro{-0.1}{4}{-0.1}{2.5}{x}{F_X(x)}
%\plotmacro{0.8}{3}{-cos{(\x*180/pi)} + 0.7}{black}
%\plotmacro{-0.1}{0.8}{0}{black}
%\plotmacro{3}{4}{1+0.69}{black}
%\draw (0,1.69) node [rotate = 90] {\tiny$|$} node [left] {$1$};
%\draw (0.8,0) node {\tiny$|$} node [below = 1mm] {$a$};
%\draw (3,0) node {\tiny$|$} node [below] {$b$};
%\end{tikzpicture}
%\end{center}
\subsection{Proof that $\boldsymbol{U = F_X(X)}$}
\noindent \textbf{Proposition:} Let $U = F_X(X)$ (where $X$ is a random variable with known distribution and cdf $F_X$). Then $U$ has a uniform distribution on $[0,1]$. \\

\noindent \textbf{Proof:} \\
Firstly, by letting $U = F_X(X)$, we have defined the outputs, $u$, of $U$, to be $u = F_X(x)$. Now, the outputs of $F_X$, like any cdf, must be between 0 and 1, so it follows that $u \in [0,1]$. \\

\noindent Let $F_U$ represent the cdf of $U$. In the following proof we restrict the doman of $F_X$ to the interval [$a,b$] on which $F_X$ is increasing. Since $F_X$ is increasing, we know that $F^{-1}_X$ is also increasing (see section \ref{increasing fact II}), which is why the direction of the inequality does not change in step 3 below
\begin{align*}
F_U(u) & = P(U < u) \\
	& = P(F_X(X) < u) \tag{since $U = F_X(x)$} \\
	& = P(X < F^{-1}_X(u)) \tag{step 3} \\
	& = F_X(F^{-1}_X(u)) \tag{by defintion of a cdf, i.e., $F_X(x) = P(X < x)$} \\
	& = u. \tag{since inverse functions undo eachother}
\end{align*}
So, for all $u \in [0,1]$ we have $F_U(u) = u$. Since $u \in [0,1]$, it also follows that $P(U~<~0)~=~0$ and that $P(U < 1) = 1$. So, we have
\[
F_U(u) = \left\{\begin{array}{ll}
0, & \; \text{if} \;\; u < 0  \smallskip \\
u, & \; \text{if} \;\; u \in [0,1]  \smallskip \\ 
1, & \; \text{if} \;\; u > 1
\end{array} \right.
\]
Which is the cdf of U($0,1$), as required.
\subsection{Proof that $\boldsymbol{X = F^{-1}_X(U)}$}
\textbf{Proposition:} Let $X = F^{-1}(U)$, where $U$ is a uniform random variable on $[0,1]$. Then $X$ has cdf $F$. \\

\noindent \textbf{Proof:} \\
Let $X =  F^{-1}(U)$, for some cdf $F$. We again restrict the domain of $F$ to the interval on which $F$ is increasing, so the inequality remains unchanged in step 3. In step 5 we use the fact that $F(x) \in [0,1]$, since $F$ is a cdf (and cdfs by definition must only have outputs between 0 and 1)
\begin{align*}
F_X(x) & = P(X < x) \\
	& = P(F^{-1}(U) < x) \tag{since $X = F^{-1}(U)$} \\
	& = P(U < F(x)) \tag{step 3} \\
	& = F_U(F(x)) \\
	& = F(x). \tag{step 5, since $F_u(u) = u$ for $u \in [0,1]$}
\end{align*}
So $X$ has cdf $F$, as required.


\section{Expected values} \label{measures of location} \label{expected value} \label{E(X)}
In section \ref{mode definition} we learned about the mode, and in section \ref{cdf strictly increasing 1} we learned about the median. Both of these quantities are known as \textbf{measures of location}: The mode tells us which position, on the $x$-axis, corresponds to the most likely output; the median tells us which position has values falling above and below it with equal (50\%) probability. Their values depend on the parameters of the model. \\

\noindent Measures of location don't tell us how `spread out' the distribution is. \\

\noindent Another widely used measure of location is the \textbf{expected value}, $\boldsymbol{E(X)}$, of a random variable, $X$. It is also known as the \textbf{mean}, and given the symbol $\boldsymbol{\mu}$. The expected value is a real number which, intuitively speaking, tells us what the average outcome of an experiement will (approximately) be, if we repeat the experiment many times. \\

\noindent In lecture 34 we will learn about the `law of large numbers,' which in most cases guarantees that, if an experiment is repeated $n$ times, then the average outcome of all the repetitions will eventually converge to the expected value, as $n$ goes to infinity. \\

\noindent Note that the `\textbf{average}' is sometimes referred to as the `\textbf{arithmetic mean}'; this is not to be confused with the `mean', which, as we have just stated, is another word for expected value. Most people will be familiar with the average: It is the sum of all results of an experiment, divided by the number of trials. \\

\noindent The expected value is defined differently for discrete and continuous distributions:

\subsubsection{$\boldsymbol{E(X)}$ when $\boldsymbol{X}$ is discrete}
For a discrete random variable, we define the expected value (i.e., the mean) to be

  	\bigskip
	
	\hfil\fbox{
	\begin{minipage}{0.35\columnwidth}{
	
	\vskip 0.04cm
	\begin{center}
	$\displaystyle{E(X) = \sum_i x_i P(X = x_i)}$
	\end{center}
	}
	\end{minipage}}
	\bigskip
	
\noindent If the number of possible values of $x$ is infinite, then this sum is infinite, in which case the definition only makes sense if the sum is convergent (an infinite sum that is not convergent will not be `equal' to anything). \\

\noindent This definition is also sometimes written as

  	\bigskip
	
	\hfil\fbox{
	\begin{minipage}{0.35\columnwidth}{
	
	\vskip 0.04cm
	\begin{center}
	$\displaystyle{E(X) = \sum_x xp(x)}$
	\end{center}
	}
	\end{minipage}}
	\bigskip

\noindent Which means the same thing (the subscript $x$ under the summation symbol means `sum over all possible $x$.' Also, recall that $p(x)$ represents a function rule which gives $P(X = x)$, see section \ref{pp(k)P(X=k)f(k)F(K)} to revise our use of this notation). \\

\subsubsection{$\boldsymbol{E(X)}$ when $\boldsymbol{X}$ is continuous}
For a continuous random variable, we define the expected value (i.e., the mean) to be

  	\bigskip
	
	\hfil\fbox{
	\begin{minipage}{0.35\columnwidth}{
	
	\vskip 0.04cm
	\begin{center}
	$\displaystyle{E(X) = \int_{-\infty}^\infty xf(x) \; dx}$
	\end{center}
	}
	\end{minipage}}
	\bigskip

\noindent Provided that this integral exists.

\subsection{Example (discrete)} \label{ex22.1}
Consider an outcome space $\Omega = \set{\omega_1,\omega_2,\omega_3,\omega_4,\omega_5}$ for which the probabilities are specified in the table below, along with the corresponding values of a random variable, $X$:
\[
\begin{array}{|c|c|c|c|c|c|} \hline
\boldsymbol{\omega}       & \omega_1   & \omega_2  & \omega_3& \omega_4& \omega_4 \\ \hline
\boldsymbol{P(\omega)}  & \frac{1}{8}&\frac{1}{2}&\frac{1}{8}&\frac{1}{16}&\frac{3}{16} \\ \hline
\boldsymbol{X(\omega)}  & -2         	     &      -1          &         1         &       2           &     0     \\ \hline
\end{array}
\]
\begin{enumerate}[(a)]
\item Find the probabilities for $Y = X^2$
\item Calculate $E(X)$
\item Calculate $E\brac{X^2}$
\item Calculate $\squac{E(X)}^2$
\end{enumerate}
\bigskip

\solution[to (a)]
As a side note, this is a good example of a non-invertible function $Y = g(X)$: We have, for example, $(-2)^2 = 4$ and $2^2 = 4$ (that is, $-2$ and 2 both map to the same element, which is the definition of a function not being invertible). \\

\noindent So, $\set{X = 2}$ and $\set{X = -2}$ are both mapped to $\set{Y = 4}$. Hence, using the values from the table,
\[
P(Y = 4) = P(X = -2) + P(X = 2) = \frac{1}{8} + \frac{1}{16} = \frac{3}{16}.
\]
Similarly, $(-1)^2 = 1^2 = 1,$ so 
\[
P(Y = 1) = P(X = -1) + P(X = 1) = \frac{1}{2} + \frac{1}{8} = \frac{5}{8}.
\]
Finally, $0^2 = 0$, so
\[
P(Y = 0) = P(X = 0) = \frac{3}{16}.
\]
We have found a corresponding value of $Y$ for each value of $X$, so we must have a value of $Y$ for every event $\omega$, so we are done. Notice that $P(Y = 0) + P(Y = 1) + P(Y = 4) = \frac{3}{16} + \frac{5}{8} + \frac{3}{16} = 1$. So, the probabilities for $Y$ sum to 1, which indicates that our calculations have been accurate (since we always know that $P(\Omega) = 1$). \\
\bigskip

\solution[to (b)]
Using the definition of $E(X)$ (discrete case) we have
\begin{align*}
E(X)  & = \sum_i x_iP(X = x_i) \\
	& = X(\omega_1)P(\omega_1) + X(\omega_2)P(\omega_2) + X(\omega_3)P(\omega_3) + X(\omega_4)P(\omega_4) + X(\omega_5)P(\omega_5) \\
	& = -2\times\frac{1}{8} + -1\times\frac{1}{2} + 1\times\frac{1}{8} + 2\times\frac{1}{16} + 0\times\frac{3}{16} \tag{from the table} \\
	& = -\frac{1}{2}.
\end{align*}
So, if we repeated the experiment many times, we would expect the average of all the results to be roughly $-\frac{1}{2}$. This makes sense because, looking at the table, we can see that, generally, higher probabilities are associated with the negative values of $X$; for example, $P(X = -1) = \frac{1}{2}$. \\
\bigskip

\solution[to (c)]
We found in the solution to (a) that $Y \in \set{0,1,4}$. So, using the definition of $E(X)$ again,
\begin{align*}
E(X^2) & = E(Y)  \\
	    & = \sum_i y_i P(Y = y_i) \\
	    & = 0\times P(Y = 0) + 1\times P(Y = 1) + 4\times P(Y = 4) \tag{since $Y \in \set{0,1,4}$} \\
	    & = 0\times\frac{3}{16} + 1\times\frac{5}{8} + 4\times\frac{3}{16} \tag{using results from (a)} \\
	    & = \frac{11}{13}. 
\end{align*}
So, \textit{the expected value has changed}, even though the events themselves and the probabilities associated with them have not changed. This is because, in this case, the expected value tells us what we should expect from $X^2$ (the square of the results), rather than from $X$. \\
\bigskip

\solution[to (d)]
This question is here to illustrate that $E(X^2) \neq \squac{E(X)}^2$. Calculating the latter is easy, since it is just the square of $E(X)$, which we found in the solution to (b). We have
\[
\squac{E(X)}^2 = \brac{-\frac{1}{2}}^2 = \frac{1}{4}.
\]


\subsection{Expected value of Poisson distribution (discrete)} \label{meanPoisson}
Recall that the probability mass function of Poi($\lambda$) is
\[
P(X = k) = \frac{\lambda^k}{k!}\e^{-\lambda}, \;\;\; \text{where} \; k \in \set{0,1,2,\ldots} 
\]
Where $X$ measures the number of some kind of occurrence, and $\lambda$ is the average number of occurrences per unit of continuum (e.g., messages in a network per second). \\

\noindent What do we predict $E(X)$ to be? We should find that after many repetitions, the average of all of our observations will approach $\lambda$. Now, using the definition of $E(X)$, but starting at $k = 1$, rather than $k = 0$, since the term with $k = 0$ will be equal to $0\times P(X = 0) = 0$,
\begin{align*}
E(X)  & = \sum_{k = 1}^\infty kP(X = k) \\
	& = \sum_{k = 1}^\infty k \frac{\lambda^k}{k!}\e^{-\lambda} \tag{pdf of Poi($\lambda$)} \\
	& = \e^{-\lambda} \sum_{k = 1}^\infty \frac{\lambda^k}{(k-1)!} \tag{linearity, and $\frac{k}{k!} = \frac{1}{(k-1)!}$ cancelling the $k$}  \\
	& =  \e^{-\lambda} \sum_{m = 0}^\infty \frac{\lambda^{m+1}}{m!} \tag{letting $m = k-1$ (so then $k = m+1$)} \\
	& =  \lambda \e^{-\lambda} \sum_{m = 0}^\infty \frac{\lambda^{m}}{m!} \tag{linearity, and $\lambda^{m+1} = \lambda^m\lambda$ (index laws)} \\
	& = \lambda \e^{-\lambda} \e^\lambda \phantom{\sum_m^\infty} \tag{power series definition of $\e^x$, section \ref{power series}}  \\
	& = \lambda. \phantom{\sum_m^\infty} \tag{since $\e^{-\lambda} \e^\lambda = \e^0 = 1$ (index laws)}
\end{align*}
So $E(X) = \lambda$, which is what we predicted.

\subsection{Expected value of Gamma distribution (continuous)} \label{meanGamma}
Recall that the probability density function of Gamma($\alpha,\lambda$) is
\[
f(x) = \frac{\lambda^\alpha}{\Gamma(\alpha)}x^{\alpha-1}\e^{-\lambda x}, \;\;\; \text{for} \; x \geq 0,
\]
and $f(x) = 0$ otherwise. Also, remember that $\Gamma(\alpha + 1) = \alpha\Gamma(\alpha)$ (from section \ref{gamma properties}). \\
We use the definition of $E(X)$ (the continuous case), except we only integrate from $0$ to $\infty$, since the integral from $-\infty$ to $0$ is equal to $0$ (as $f(x) = 0$ when $x < 0$),
\begin{align*}
E(X)  & = \int_0^\infty x \frac{\lambda^\alpha}{\Gamma(\alpha)}x^{\alpha-1}\e^{-\lambda x} \; dx \\
	& = \int_0^\infty \frac{\lambda^\alpha}{\Gamma(\alpha)}x^{\alpha}\e^{-\lambda x} \; dx \tag{since $xx^{\alpha - 1} = x^\alpha$} \\
	& = \frac{\alpha}{\lambda}  \int_0^\infty \frac{\lambda}{\alpha} \frac{\lambda^\alpha}{\Gamma(\alpha)}x^{\alpha}\e^{-\lambda x} \; dx \tag{linearity, and multiplying by $\frac{\alpha}{\lambda}\frac{\lambda}{\alpha} = 1$} \\
	& = \frac{\alpha}{\lambda}  \int_0^\infty \frac{\lambda^{\alpha+1}}{\Gamma(\alpha + 1)}x^{\alpha}\e^{-\lambda x} \; dx. \tag{since $\alpha\Gamma(\alpha) = \Gamma(\alpha + 1)$ and $\lambda \lambda^\alpha = \lambda^{\alpha+1}$}
\end{align*}
On the right hand side we have $\frac{\alpha}{\lambda}$ multiplied by the integral of the pdf of Gamma($\alpha+1,\lambda$). Since this is a pdf, we know that the integral over all possible values for $x$ must be equal to 1, by probability measure axiom 1. Hence
\[
E(X) = \frac{\alpha}{\lambda}.
\]
Note that the mode (most probable value) of Gamma($\alpha,\lambda$) is $\frac{\alpha-1}{\lambda}$. This is a good example of a case where the mode is not equal to the mean. You can easily calculate the mode using calculus (i.e., by finding $f'(x)$ then solving for $x$ in $f'(x) = 0$). \\

\subsubsection{Mean of exponential distribution}
\noindent Let $X \sim$ Exponential($\lambda$). The exponential distribution is the special case of Gamma($\alpha,\lambda$), with $\alpha = 1$ (we showed this at the end of section \ref{gamma distribution1}). Hence we can find the mean of the exponential distribution:
\[
\text{For Exponential($\lambda$),} \;\; E(X) = \frac{\alpha}{\lambda}, \; \text{with $\alpha = 1$} \;\; \text{i.e.,} \;\;\;\; E(X) = \frac{1}{\lambda}.
\]

\section{Expected value of $\boldsymbol{g(X)}$} \label{def E(g(X))}
In example \ref{ex22.1} we manually calculated $E(g(X))$ for the case $g(X) = X^2$, by first finding all of the probabilities for $Y = X^2$. In general, it will be easier to use the following definitions of $E(g(X))$: \\

\noindent\textbf{When $\boldsymbol{X}$ is discrete}

  	\bigskip
	
	\hfil\fbox{
	\begin{minipage}{0.4\columnwidth}{
	
	\vskip 0.04cm
	\begin{center}
	$\displaystyle{E(g(X)) = \sum_i g(x_i) P(X = x_i)}$
	\end{center}
	}
	\end{minipage}}
	\bigskip

\noindent \textbf{When $\boldsymbol{X}$ is continuous}

  	\bigskip
	
	\hfil\fbox{
	\begin{minipage}{0.35\columnwidth}{
	
	\vskip 0.04cm
	\begin{center}
	$\displaystyle{E(g(X)) = \int_{-\infty}^\infty g(x)f(x) \; dx}$
	\end{center}
	}
	\end{minipage}}
	\bigskip

\noindent Note that, while we refer to these as `definitions' in this subject, they are actually a result of a theorem called the Law of the Unconscious Statistician, which will not be covered in this course. 

\subsection{Important properties of expected value} \label{sec22.01}
In general, it is unfortunately often true that
\[
E(g(X)) \neq g(E(X)).
\]
We saw this in example \ref{ex22.1}, parts (c) and (d), where we learned that
\[
E(X^2) \neq \squac{E(X)}^2 
\]
Which has $g(X) = X^2$. \\

\noindent The following simple properties, however, are true in general, and will be very useful to us. \\

\noindent We can prove each of these properties using the above definition of $E(g(X))$:
\begin{enumerate}
\item $E(c) = c$ for all $c \in \R$
\item Linearity properties:
\begin{itemize}
\item $E(kX) = kE(X)$
\item $E(a + X) = E(a) + E(X)$
\end{itemize}
\item $E(kX + a) = kE(X) + a$ (this is just a consequence of the above properties)
\end{enumerate}

\noindent Notice that, for the continuous case, the expected value is defined in terms of an integral. For the discrete case, it is defined in terms of a sum. Since integrals and sums are both linear, we should not be too surprised that the expected value is also linear, but be warned: Thinking that something is true, is never a good substitute for proving it.

\subsubsection{Property (1)}
\noindent \textbf{Proof (discrete case):}
Let $X$ be a discrete random variable. Let $c \in \R$. Then
\begin{align*}
E(c)  & = \sum_i c P(X = x_i) \tag{definition of $E(g(X))$ with $g(x) = c$} \\
	& = cP(X = x_1) + cP(X = x_2) + \ldots \\
	& = c\brac{P(X = x_1) + P(X = x_2) + \ldots} \tag{factoring out $c$} \\
	& = c. \tag{since $P(\Omega) = 1$}
\end{align*}

\noindent The proof of the continuous case is left as an exercise.

\subsubsection{Property (2): Linearity}

\noindent \textbf{Proof (continuous case):}
Let $X$ be a continous random variable. Let $k \in \R$. Then, using the definition of $E(g(X))$, with $g(x) = kx$ we have
\[
E(kX) = \int_{-\infty}^\infty kxf(x) \; dx = k\int_{-\infty}^\infty xf(x) \; dx = kE(X) \tag{using linearity of integrals}
\]
Now, let $a \in \R$, then using linearity of integrals again we get
\[
E(a+X) = \int_{-\infty}^\infty (a+x)f(x) \; dx = \int_{-\infty}^\infty af(x) \; dx + \int_{-\infty}^\infty xf(x) \; dx = E(a) + E(X).
\]
Hence $E(kX) = kE(X)$  and $E(a + X) = E(a) + E(X)$, so expected value is linear. \\

\noindent Proving the discrete case is left as an exercise.


\subsection{Example (mean of normal distribution)} \label{meanNormal}
Let $Z \sim$ N($0,1$). By the definition of $E(Z)$ we have
\begin{align*}
E(Z)  & = \frac{1}{\sqrt{2\pi}} \int_{-\infty}^\infty z\e^{-z^2/2} \; dz \tag{using pdf of N($0,1$)} \\
	& = 0.
\end{align*}
The reason this integral evaluates to $0$ is because $z$ in an odd function, and $\e^{-z^2/2}$ is an even function: The product of an odd function with an even function is always an odd function, and the integral over a symmetric interval of an odd function is always equal to 0 (see section \ref{symmetric interval}). \\

\noindent Now, let $X \sim$ N($\mu,\sigma^2$). Recall that $X$ can be transformed into $Z$ (standard normal) by letting
\[
Z = \frac{X - \mu}{\sigma} \tag{section \ref{sec21.2}}
\]
\noindent Using this, and the fact that $E(Z) = 0$, calculate $E(X)$. \\

\solution
Solving for $X$ in the equation above gives, $X = \sigma Z + \mu$. Hence,
\[
E(X) = E(\sigma Z + \mu) = \sigma E(Z) + \mu = \mu. \tag{since $E(Z) = 0$}
\]
So the mean of the normal distribution is $\mu$.

%\noindent For example, if a standard, 6 sided, balanced die, is rolled many times, then eventually the average of all of the outcomes of the rolls will approach $3.5$, which (we will soon see) is the expected value.





